A tangent ray through $(0,16)$ has line $y=mx+16$. Its distance from $(6,10)$ must be 2, giving $(6m+6)^2=4(m^2+1)$, or $4m^2+9m+4=0$. The two slopes have product 1 and difference $\dfrac{\sqrt{17}}{4}$. Their x-intercepts are $- \dfrac{16}{m_1} , - \dfrac{16}{m_2}$, so the shadow length is $16 | m_1 - \dfrac{m_2 |}{| m_1 m_2 |} = 4 \sqrt{17}$. Thus $\boxed{m+n=21}$.
\begin{asy}
import graph;
usepackage("amsmath");
import geometry;
size(7cm);
pair S=(0,16),O=(6,10);draw((-1,0)--(28,0));draw(circle(O,2));real m1=(-9+sqrt(17))/8,m2=(-9-sqrt(17))/8;draw((-16/m1,0)--S--(-16/m2,0));label("$S(0,16)$",S,(0,1));label("$O(6,10)$",O,(1,0));label("shadow",(15,0),(0,-1));
\end{asy}

